% ID: SEL-D0022-OA-E03
% ESTADO: validada. S023, revisada con los criterios oficiales.
Sean $F$ y $W$ usar Facebook y WhatsApp, respectivamente.

\textbf{a)} $F$ se descompone en $F\cap W$ y $F\cap W^c$,
así que $P(F\cap W)=0.60-0.04=0.56$, es decir, $\boxed{56\%}$.

\textbf{b)} De los que usan WhatsApp, quitamos los que usan ambas.
$P(W\cap F^c)=0.80-0.56=0.24$, es decir, $\boxed{24\%}$.

\textbf{c)} Ahora el grupo de referencia son los usuarios de WhatsApp.
$P(F\mid W)=0.56/0.80=0.70$, es decir, $\boxed{70\%}$.

\textbf{d)} Si fueran independientes, la intersección tendría
probabilidad $0.60\cdot0.80=0.48$. Como la observada es $0.56$,
los sucesos son dependientes.
